\documentclass{report}

\usepackage{biblatex}
\usepackage{hyperref}
\usepackage{minted}
\usepackage[bottom]{footmisc}
\usepackage{amsmath}
\usepackage{enumitem}


% \addbibresource{bibliography.bib}

\setlength{\parskip}{\bigskipamount}

\begin{document}

\tableofcontents

\chapter{Introduction}
In this report, a geometric coverage problem is presented and solved using various combinatorial optimization techniques. The problem consists of a set of rectangles and their respective vertices. The primary objective is to deploy a set of guards strictly on these vertices such that every rectangle is covered. A rectangle is considered covered if at least one of its incident vertices possesses a guard. Consequently, the initial problem is to determine a vertex assignment that minimizes the total number of guards required.

Subsequently, an extension to this baseline problem is explored. In practical contexts, guards may represent resources that require differentiation, preventing interference or scheduling conflicts. Therefore, a secondary objective is introduced, which consists of assigning a color to each placed guard such that no rectangle is covered by more than one guard of the same color, while minimizing the total number of colours utilized. This formulation effectively transforms the problem into a multi-objective optimization problem.

The following chapters establish the mathematical formulations using Greedy Approaches, Integer Programming, and Constraint Programming. Then, the implementation of these approaches using Google OR-Tools and an independent search method maintaining arc-consistency with AC-3 is described. Finally, the multi-objective extension is formulated and solved using a sequential optimization method.

\section{Objectives}
The primary objectives of this report are defined as follows:

\begin{enumerate}
    \item To evaluate Greedy Approaches as a preliminary method for the guard placement problem.
    \item To establish rigorous mathematical formulations for the problem utilizing both Integer Programming and Constraint Programming models.
    \item To implement the proposed mathematical models utilizing the Google OR-Tools suite.
    \item To develop and evaluate a custom search algorithm employing AC-3 for information filtering and backtracking to solve the Constraint Optimization Problem.
    \item To formulate and solve a multi-objective extension of the problem, focusing on the minimization of the colours assigned to the minimum number of guards found in the initial models.
\end{enumerate}


\chapter{Greedy Approaches}

\section{Analysis}


\chapter{Integer Programming}
In this chapter, an Integer programming model is presented. The initial step in this process is to establish the mathematical formulation of the problem, which includes the definition of the decision variables, constraints, and objective function.
Then, the implementation of this approach using Google OR-Tools is described, followed by an interpretation of the results.
\section{Mathematical Model}
Let $m$ be the number of rectangles to cover, $n$ the number of vertices and $p_{ij}$ represented as follows:

\[
p_{ij} =
\begin{cases}
1, & \text{if rectangle } j \text{ is incident to vertex } i \\
0, & \text{otherwise}
\end{cases}
\]

Decision Variables: $v_{i}$ indicates wheter vertex $i$ has a guard or not.

The objective is to minimize the total number of guards:

\[\sum_{i = 1}^{n} v_{i}\]

subject to the following constraints:

\[
\begin{cases}
	\sum_{i = 1}^{n} v_{i} \times p_{ij} \ge 1, \quad \forall j = 1,2,\dots,m \\
	v_{i} \in \{0,1\}, \quad \forall i = 1,2,\dots,n                          \\
\end{cases}
\]

The objective of this constraint is to ensure that each rectangle is covered by ensuring that at least one of its vertices possesses a guard.

\section{OR-Tools}
The implementation of the Integer Programming model utilizes the \texttt{MPSolver} interface from Google OR-Tools, specifically configured to use the SCIP (Solving Constraint Integer Programs) backend.

The implementation translates the mathematical formulation directly into code:
\begin{itemize}
    \item \textbf{Variables:} Using \texttt{solver->MakeIntVar(0.0, 1.0, name)}, boolean decision variables $v_i$ are instantiated for each candidate vertex in the problem space.
    \item \textbf{Constraints:} For each rectangle, a row constraint is created via \texttt{solver->MakeRowConstraint(1.0, solver->infinity())}. The variables corresponding to the vertices incident to that rectangle are assigned a coefficient of 1. This enforces the condition $\sum_{i=1}^{n} v_i \times p_{ij} \ge 1$.
    \item \textbf{Objective Function:} The objective function is defined by setting the coefficient of every vertex variable to 1 and invoking \texttt{objective->SetMinimization()} to ensure the solver finds the minimum set of guards.
\end{itemize}

Additionally, to improve solver performance by reducing the search space, a theoretical lower bound for the number of guards needed (e.g., $\lceil m / 3 \rceil$, where $m$ is the number of rectangles to be covered) is calculated and added as a global row constraint. This prevents the solver from evaluating unfeasible sub-optimal branches.

\chapter{Constraint Programming}
% update
In this chapter, a Constraint programming model is presented. The initial step in this process is to establish the Constraint Optimization Problem (COP). Then, we discuss how to solve the COP using two different methods: maintaining arc-consistency with AC-3 and using OR-Tools.

\section{Constraint Optimization Problem}
A Constraint Optimization Problem is a tuple $(X, D, C)$ where:
\begin{itemize}
    \item $X = \{x_1, x_2, \dots, x_n\}$ is the set of variables;
    \item $D = \{d_1, d_2, \dots, d_n\}$ is the set of domains, where $d_i$ is a finite set of possible values for $x_i$;
    \item $C = \{c_1, c_2, \dots, c_m\}$ is the set of constraints;
    \item $f$ is the optimization function.
\end{itemize}

The following COP can be used to model this problem:

\begin{itemize}
  \item $X = \{v_1, v_2, \dots, v_{n}\}$: set of vertices;
    \item $D = \{d_1, d_2, \dots, d_{n}\}: d_1 = d_2 = \dots = d_{n} = \{0,1\}$;
    \item $C$: \(\sum_{i = 1}^n v_i \times p_{ij} \ge 1 \in C, \quad \forall j = 1,2,\dots,m \);
    \item $f$ = \(\sum_{i = 1}^n v_i\). 
\end{itemize}

The model under discussion is essentially analogous to the model for integer programming, but expressed in the form of a COP.
\section{AC-3}
\section{OR-Tools}
To implement the Constraint Optimization Problem, the OR-Tools CP-SAT solver is employed via the \texttt{CpModelBuilder} API. The CP-SAT solver utilizes satisfiability (SAT) methods under the hood to efficiently explore the domain.

The model construction follows a two-phase approach based on the implementation:
\begin{enumerate}
    \item \textbf{Guard Placement (Phase 1):}
    \begin{itemize}
        \item \textbf{Variables:} The decision variables $v_i$ are initialized using \texttt{cp\_model.NewIntVar(domain)}, strictly bounded to a domain of $\{0, 1\}$.
        \item \textbf{Constraints:} The coverage constraints are constructed using linear expressions (\texttt{LinearExpr}). For each rectangle, the boolean variables of its incident vertices are summed, and the constraint is strictly enforced using \texttt{cp\_model.AddGreaterOrEqual(expr, 1)}.
        \item \textbf{Objective:} The objective is defined as minimizing the sum of all $v_i$ variables via \texttt{cp\_model.Minimize(expr)}.
    \end{itemize}

    \item \textbf{Guard Coloring (Phase 2):}
    Once the first phase determines the optimal minimum number of guards, a secondary solver instance is invoked to solve an assignment problem (e.g., assigning distinct frequencies or shifts), effectively minimizing the chromatic number of the placed guards.
    \begin{itemize}
        \item \textbf{Variables:} The model introduces a boolean array $c_j$ indicating whether a color $j$ is utilized, and a 2D matrix $u_{i,j}$ mapping vertex $i$ to color $j$.
        \item \textbf{Constraints:} The constraints ensure that the total number of placed guards strictly equals the optimal solution found in Phase 1 (\texttt{AddEquality(expr, sol)}), and bounds the color assignment variables to logically link the mapping $u_{i,j}$ with the active colors $c_j$.
        \item \textbf{Objective:} The goal transitions to minimizing the sum of the color variables $c_j$.
    \end{itemize}
\end{enumerate}

The models are compiled and solved using \texttt{Solve(cp\_model.Build())}. The solver returns a \texttt{CpSolverResponse}, which yields both the optimal geometric configuration of guards and their minimum chromatic assignment.


\chapter{Extensions}

\chapter{Conclusion}

% \printbibliography

\end{document}
